\section{Buildprozess}
\label{sec:buildprozess}

Das gesamte System wird über Docker und Docker Compose gebaut und betrieben.
Jede Komponente besitzt ein eigenes Dockerfile, das das zugehörige Image definiert.
Docker Compose orchestriert alle Komponenten gemeinsam und kümmert sich um
Netzwerkkonfiguration, Startreihenfolge und Health-Checks.

\subsection{Verzeichnisstruktur und Dockerfiles}

Die Dockerfiles liegen direkt bei den jeweiligen Komponenten:

\begin{lstlisting}[language={},caption={Verzeichnisstruktur der Dockerfiles}]
VerteilteSysteme-TaskGrid-/
  Dockerfile                    -- Dispatcher
  src/
    namensdienst/
      Dockerfile                -- Namensdienst
    worker/
      Dockerfile                -- Alle Worker-Typen
      entrypoint.sh             -- Worker-Startskript (setzt Env-Vars)
    client/
      Dockerfile                -- Client
  docker-compose.yml            -- Orchestrierung aller Dienste
\end{lstlisting}

\subsection{Dispatcher-Image}

Das Dispatcher-Dockerfile liegt im Wurzelverzeichnis des Projekts.
Es basiert auf \texttt{python:3.11-slim} und installiert alle Python-Abhängigkeiten
aus \texttt{requirements.txt}. Die Proto-Stubs (\texttt{proto/}) sowie der
gesamte Quellcode werden in den Container kopiert.

\begin{lstlisting}[language=bash,caption={Image manuell bauen (optional)}]
docker build -t taskgrid-dispatcher .
\end{lstlisting}

\subsection{Namensdienst-Image}

Der Namensdienst hat ein eigenes Dockerfile unter \texttt{src/namensdienst/Dockerfile}.
Es setzt den Python-Pfad so, dass sowohl \texttt{src/} als auch \texttt{proto/}
importierbar sind. Die Umgebungsvariable \texttt{PYTHONUNBUFFERED=1} sorgt dafür,
dass Log-Ausgaben sofort in \texttt{docker logs} erscheinen und nicht gepuffert werden.

\subsection{Worker-Image und Entrypoint}

Alle Worker-Typen teilen sich ein einziges Dockerfile unter \texttt{src/worker/Dockerfile}.
Was einen \texttt{worker-sum} von einem \texttt{worker-reverse} unterscheidet,
ist ausschließlich die Umgebungsvariable \texttt{WORKER\_TYPE}, die beim Start
in der \texttt{docker-compose.yml} gesetzt wird.

Das Startskript \texttt{entrypoint.sh} übernimmt dabei eine wichtige Aufgabe:
Es übersetzt die generische Variable \texttt{WORKER\_TYPE} in die für den Worker
nötigen spezifischen Variablen und generiert eine eindeutige Worker-ID aus Typ
und Hostname des Containers:

\begin{lstlisting}[language=bash,caption={Auszug aus entrypoint.sh}]
export TASK_TYPES="${WORKER_TYPE}"
export NAMING_SERVICE_ADDRESS="${NAMENSDIENST_ADDRESS}"
export WORKER_ID="worker-${WORKER_TYPE}-$(hostname)"
exec python src/worker/worker.py
\end{lstlisting}

Durch den hostname-basierten Anteil der Worker-ID sind alle Instanzen desselben
Worker-Typs eindeutig identifizierbar. Auch wenn drei \texttt{worker-sum}-Container
gleichzeitig laufen.

\subsection{docker-compose.yml}

Die \texttt{docker-compose.yml} im Wurzelverzeichnis definiert alle acht Dienste
des Systems und legt fest, wie sie miteinander kommunizieren:

\begin{lstlisting}[language=yaml,caption={Dienste in docker-compose.yml}]
services:
  namensdienst     # Service Discovery (startet zuerst)
  dispatcher       # Aufgabenverwaltung und Dispatch
  worker-sum       # Worker fuer Typ "sum"   (skalierbar)
  worker-reverse   # Worker fuer Typ "reverse"
  worker-hash      # Worker fuer Typ "hash"
  worker-upper     # Worker fuer Typ "upper"
  worker-wait      # Worker fuer Typ "wait"
  client           # Demo-Client (One-Shot)
\end{lstlisting}

\subsubsection*{Startreihenfolge und Health-Checks}

Damit die Komponenten in der richtigen Reihenfolge hochfahren, ist jeder Dienst
mit einem \texttt{healthcheck} ausgestattet. Der Health-Check öffnet eine
TCP-Verbindung auf den gRPC-Port. Das ist deutlich zuverlässiger als nur zu
prüfen, ob der Prozess läuft, weil damit sichergestellt ist, dass der gRPC-Server
tatsächlich Verbindungen annimmt.

\begin{lstlisting}[language=yaml,caption={Health-Check (Beispiel Namensdienst)}]
healthcheck:
  test: ["CMD", "python", "-c",
         "import socket; s=socket.socket(); s.settimeout(2);
          s.connect(('localhost',50052)); s.close()"]
  interval: 5s
  timeout: 3s
  retries: 12
  start_period: 5s
\end{lstlisting}

Die \texttt{depends\_on}-Bedingungen nutzen \texttt{condition: service\_healthy},
sodass Docker Compose erst dann mit dem nächsten Dienst weitermacht, wenn der
vorherige seinen Health-Check bestanden hat.

\subsubsection*{Netzwerk}

Alle Dienste sind in einem gemeinsamen Bridge-Netzwerk namens \texttt{taskgrid}
verbunden. Innerhalb dieses Netzwerks können Dienste sich gegenseitig über ihren
Service-Namen erreichen \texttt{Dispatcher} löst also zu der Container-IP des
Dispatchers auf, \texttt{namensdienst} entsprechend zum Namensdienst-Container.
Nach außen sind nur der Dispatcher (Ports 50051 und 8080) und der Namensdienst
(Port 50052) exponiert.

\subsubsection*{Volumes}

Laut Aufgabenstellung dürfen gemeinsame Volumes nicht für die Task- oder
Ergebnis-Übergabe genutzt werden. Das einzige Volume im System ist deshalb
\texttt{logs}, das die strukturierten Log-Dateien des Dispatchers und Clients
über Container-Neustarts hinaus persistiert.

\subsubsection*{Konfiguration über Umgebungsvariablen}

Alle Dienste werden ausschließlich über Umgebungsvariablen konfiguriert,
keine Konfigurationsdatei wird gemountet. Das macht es möglich, das Verhalten
einzelner Dienste zu ändern, ohne Images neu zu bauen. Zum Beispiel den
Heartbeat-Timeout des Namensdiensts für Tests auf einen niedrigeren Wert zu setzen:

\begin{lstlisting}[language=yaml,caption={Konfiguration des Namensdiensts (Ausschnitt)}]
environment:
  - NAMESERVICE_PORT=50052
  - HEARTBEAT_TIMEOUT_SECONDS=15
  - HEARTBEAT_CHECK_INTERVAL_SECONDS=5
  - PYTHONUNBUFFERED=1
\end{lstlisting}

\subsection{Skalierung einzelner Worker-Typen}

Da die Worker-Container keinen \texttt{container\_name} haben, kann Docker Compose
mehrere Instanzen desselben Dienstes starten. Dabei vergibt Docker automatisch
eindeutige Container-Namen. Der interne DNS-Name (\texttt{worker-sum}) bleibt
für alle Instanzen gleich. Docker Compose löst ihn per Round-Robin auf alle
laufenden Instanzen auf:

\begin{lstlisting}[language=bash,caption={Worker-Typ skalieren}]
docker compose up --build -d --scale worker-sum=3
\end{lstlisting}

Jede der drei \texttt{worker-sum}-Instanzen registriert sich beim Namensdienst
mit einer eigenen Worker-ID. Der Dispatcher bekommt bei \texttt{LookupWorker}
alle drei zurück und verteilt Tasks gleichmäßig auf sie. Ohne dass am
Dispatcher oder Namensdienst irgendetwas konfiguriert werden muss.

\subsection{Reproducibility}

Damit der Build auf verschiedenen Rechnern reproduzierbar ist, werden die
Python-Abhängigkeiten über \texttt{requirements.txt}-Dateien gepinnt.
Die Proto-Stubs werden nicht dynamisch generiert, sondern als fertige
\texttt{*\_pb2.py}-Dateien im Repository eingecheckt. Das verhindert, dass
unterschiedliche \texttt{grpcio-tools}-Versionen zu inkompatiblen Stubs führen.
